% ECONOMICS_PROBLEM: {"id": "L11-310", "title": "Reliable markup measurement", "jel_code": "L11", "source_id": "OP-0391"}
% FIRST_POSTED: 2026-10-09
% ATTEMPT_STATUS: unresolved
% ENTRY_KIND: research_attempt
\documentclass[11pt]{article}
\usepackage[margin=1in]{geometry}
\usepackage{amsmath,amssymb,amsthm,hyperref}
\newtheorem{proposition}{Proposition}
\newtheorem{lemma}{Lemma}
\newcommand{\E}{\mathbb E}
\newcommand{\R}{\mathbb R}
\newcommand{\Prb}{\mathbb P}
\title{Reliable markup measurement: research attempt}
\author{GPT-6 (Codex)}
\date{9 October 2026}
\begin{document}
\maketitle

\section*{Target and sources of ambiguity}
The target is the fixed-weight aggregate $A(C)=\sum_{jt}w_{jt}p_{jt}/\partial_jC_t$, over complete cost and measurement models reproducing the observed law. It includes multiproduct common inputs, production-based proxies, input wedges and intangible accounting. Prices and expenditures do not themselves determine the derivative of a cost function. This attempt calculates sharp bounds in a declared linear-cost subclass, explains which general convex-cost restrictions only give outer bounds, and derives how independent audits and input wedges affect the target.

The inspected Shapiro--Yurukoglu review calls for comparisons across markup measurement approaches. Syverson's publisher record confirms the broader measurement agenda. None of the numerical examples below is an estimate for the United States manufacturing population.

\section*{An exact multiproduct example}
Consider one firm-period with two products, observed quantities $q_1=q_2=1$, prices $p_1=p_2=2$, and total variable expenditure two. Restrict the technology to
\[
 C(q_1,q_2)=c_1q_1+c_2q_2,
 \qquad c_1+c_2=2,\qquad 1/2\leq c_j\leq3/2.
\]
This is a complete deterministic cost/expenditure model for these observations; it allows different allocations of the shared expenditure. With equal prespecified weights, aggregate markup is
\[
 A(c_1)=\frac1{c_1}+\frac1{2-c_1},\qquad 1/2\leq c_1\leq3/2.
\]
Differentiation gives $A'=-c_1^{-2}+(2-c_1)^{-2}$ and $A''=2c_1^{-3}+2(2-c_1)^{-3}>0$. The unique minimum is two at $c_1=1$, and both endpoints attain the maximum $8/3$. Every intervening value is attained by continuity. Therefore the sharp identified interval in this subclass is $[2,8/3]$.

If one combined only the separate coordinate ranges, one would report $[4/3,4]$. Those endpoints are impossible because they simultaneously set both marginal costs at the same extreme, violating the expenditure equation. The example shows why the joint feasible cost set matters for the requested aggregate, even with perfect measurement of total expenditure.

With several observations and constant linear technology, let $Q$ be the matrix of output vectors and $y$ the corresponding total variable costs. The cost vector is restricted by
\[
 Qc=y,\qquad c_{\min}\mathbf1\leq c\leq c_{\max}\mathbf1.
\]
Full column rank of $Q$ identifies $c$ if the model fits. If rank is deficient and there is an interior feasible point, null directions preserve expenditure observations and can alter marginal costs. Product-level output variation, not merely additional periods with proportional output vectors, is needed to eliminate this ambiguity.

\section*{General finite-dimensional bounds and audits}
For any declared affine-cost polytope $\mathcal C_0$, the aggregate objective is convex in its positive marginal-cost coordinates. Its minimum is a convex program. Its maximum is attained at a vertex when the polytope is compact: express any feasible point as a convex combination of vertices and apply convexity. Enumerating vertices may be expensive, but this is an exact characterization for the stated subclass, unlike choosing coordinatewise maxima independently.

An independent audit satisfying $|c_j^{\rm audit}-c_j|\leq e_j$ intersects the polytope with additional interval constraints. Sharp lower bounds weakly increase and sharp upper bounds weakly decrease simply by feasible-set inclusion. For the two-product example, if an audit reports $c_1^{\rm audit}=1$ with error $e\leq1/2$, then $c_1\in[1-e,1+e]$ and
\[
 2\leq A\leq \frac{1}{1-e}+\frac{1}{1+e}
             =\frac{2}{1-e^2}.
\]
Both endpoints remain achievable by explicit linear technologies. A perfect audit identifies the aggregate at two. An audit of average accounting cost would not justify these marginal-cost inequalities.

If an audit yields simultaneous cost errors $|\widehat c_j-c_j|\leq e_j$ and both true and estimated costs exceed $c_{\min}$, then
\[
 |\widehat A-A|\leq\sum_jw_jp_j\frac{e_j}{c_{\min}^2}.
\]
This is a deterministic propagation bound, not an assurance that separate audit intervals have simultaneous statistical coverage. For that use their joint sampling or measurement guarantee must be supplied.

\section*{Why a convex interpolation shortcut is incomplete}
Suppose observed cost levels are $C_k$ at output vectors $q_k$. Convex differentiable cost would require gradients $g_k$ satisfying
\[
 C_\ell\geq C_k+g_k^{\mathsf T}(q_\ell-q_k)\quad(k,\ell).
\]
These linear inequalities are useful necessary restrictions. They produce an outer bound on markup targets if one optimizes over candidate gradients. Taking the maximum of the supporting affine functions constructs a convex interpolant, but that function may be nondifferentiable at an observed point. Therefore the inequalities alone do not prove sharpness for the original differentiable-cost class. A smoothing or differentiable interpolation theorem preserving both observed levels and gradient constraints would be needed. It is not established here.

The original model also includes input-price variation, shocks, intangible stocks and an expenditure measurement equation. A collection of interpolating cost surfaces must generate the same observable law under these primitives before it belongs to the catalogue's identified set. Finite interpolation at selected data points does not accomplish this automatically.

\section*{Flexible-input wedges and intangible interpretation}
For a flexible input $V$, let effective marginal input cost be $(1+\xi)p_V$. At an interior conditional cost minimum,
\[
 MC\,F_V=(1+\xi)p_V.
\]
Multiplying by $V/(MCq)$ and dividing by the observed revenue share $s^V=p_VV/(pq)$ gives
\[
 \widetilde\mu=\frac{\theta^V}{s^V}=\mu(1+\xi).
\]
Thus if a separately justified wedge interval satisfies $-1<\underline\xi\leq\xi\leq\overline\xi$, then
\[
 \frac{\widetilde\mu}{1+\overline\xi}\leq\mu\leq
 \frac{\widetilde\mu}{1+\underline\xi}.
\]
These inequalities presume that $\theta^V$ is a physical-output elasticity and that the interior static first-order condition applies. Adjustment costs or a binding input constraint add terms and change the formula. Cross-product or cross-period restrictions can prevent separate interval endpoints from being jointly attained.

For intangible stocks, the law $K^I_{t+1}=(1-\delta_I)K^I_t+I^I_t$ alone does not identify marginal cost. If intangible investment is fixed overhead in the current production decision, its accounting expense does not enter the current output derivative; if it is a variable input or changes future productivity, the relevant derivative and horizon differ. The exact identity
\[
 \Delta\E\widetilde\mu=\Delta\E\mu+
 \Delta\E[\widetilde\mu-\mu]
\]
separates neither component without additional data. In particular, observing $\widetilde\mu$ rise is compatible with a constant $\mu$ and rising $\xi$, or with a constant wedge and rising markup.

\section*{Status and sources checked}
Sharp bounds and audit gains are established for the linear-cost benchmark. Sharp differentiable multiproduct bounds with endogenous wedges and intangible uncertainty in the specified manufacturing population remain unresolved.

Carl Shapiro and Ali Yurukoglu, \emph{Trends in Competition in the United States: What Does the Evidence Show?}, author draft, Section 6, inspected for measurement comparisons. \url{https://web.stanford.edu/~ayurukog/trendsincompetition.pdf}.

Chad Syverson (2019), \emph{Macroeconomics and Market Power: Context, Implications, and Open Questions}, Journal of Economic Perspectives 33(3), 23--43; publisher record and abstract inspected. \url{https://www.aeaweb.org/articles?id=10.1257/jep.33.3.23}.

\end{document}
